% Source: Wald GR, physical PDF p. 244 (printed p. 231), Figure 9.4.
% Coverage: A point conjugate to a spacelike hypersurface.
% Figures/tables/footnotes: Figure 9.4.
% Status: source-reviewed and compile-clean.
% Uncertainties: none.

\begingroup
\renewcommand{\thefigure}{9.4}
\begin{figure}[tb]
  \centering
  \begin{tikzpicture}[x=.88cm,y=.88cm,>=Stealth,thick]
    \draw (-2.2,.15) .. controls (-1.35,.72) and (-.35,.22) .. (.4,.36)
      .. controls (1.1,.47) and (1.65,.1) .. (2.55,.48);
    \node[right] at (1.75,.32) {\(\Sigma\)};
    \coordinate (p) at (-.1,3.15);
    \draw (-.7,.5) .. controls (-.53,1.35) and (-.2,2.28) .. (p);
    \draw (-.28,.42) .. controls (-.23,1.22) and (-.11,2.3) .. (p);
    \draw[->] (-.51,1.58) -- (-.29,1.58);
    \node[right] at (.22,1.58) {\(\eta^a\)};
    \draw[->] (-1.18,2.12) to[out=18,in=150] (-.58,2.27);
    \node[left] at (-1.12,2.16) {\(\gamma\)};
    \fill (p) circle (2pt) node[above left] {\(p\)};
    % Reserve trim clearance above the artwork when this float lands at page top.
    \path[draw=none] (0,4.45) circle[radius=.1pt];
  \end{tikzpicture}
  \caption{A spacetime diagram depicting a point \(p\) conjugate to the hypersurface
  \(\Sigma\) along the geodesic \(\gamma\). (The two geodesics shown are supposed to
  be ``infinitesimally nearby.'')}
  \label{fig:9.4}
\end{figure}
\endgroup
